\section{Introduction}
\label{sec:introduction}

Retail pricing ranks among the most consequential sequential decision problems in operations research and marketing. Every price change triggers a cascade of demand responses: consumers substitute across SKUs, stockpile during promotions, and competitors react. These dynamics unfold over weeks and quarters, with effects that propagate through substitution networks and depend on store demographics, seasonality, and latent competitive context. Despite decades of empirical work on price elasticity and demand estimation \citep{hoch1995determinants, nevo2001measuring, montgomery1997estimating}, the systems that support actual pricing decisions today fall into two camps. Hand-crafted analytical simulators encode structural assumptions about demand and competition, but their fidelity to real-world dynamics is limited by the tractability of closed-form solutions. Model-free reinforcement learning methods \citep{chen2022dynamic, ban2021personalized} learn policies directly from data but forego an explicit dynamics model, making counterfactual reasoning and imagination-based planning impossible. No learned world model exists for economic environments.

This gap can be formalized using a simple taxonomy. Consider the space of sequential decision environments along two axes: the domain (physical versus economic) and the nature of the dynamics model (learned versus hand-crafted). In the physical domain, learned world models have achieved remarkable success. The Dreamer family \citep{hafner2020dreamerv1, hafner2021dreamerv2, hafner2023dreamerv3} pioneered recurrent state-space models (RSSMs) for latent imagination and policy optimization. MuZero \citep{schrittwieser2020muzero} demonstrated that planning with a learned model suffices for mastering board games and Atari. More recently, IRIS \citep{zhang2023iris} and TD-MPC2 \citep{hansen2024tdmpc2} have extended these ideas to continuous control and high-dimensional observations. These methods rely on hand-crafted physical simulators such as MuJoCo, Isaac Gym, or CARLA for evaluation and transfer. In the economic domain, by contrast, the landscape is dominated by hand-crafted models: agent-based simulators like ABIDES, structural demand models in the spirit of \citet{berry1995automobile} and \citet{nevo2001measuring}, and the AI Economist \citep{zheng2020aieconomist}, which uses a hand-designed economic environment for tax policy experiments. The quadrant of learned models applied to economic domains has remained empty. \dreamprice{} fills this gap.

\begin{table}[t]
\centering
\caption{Taxonomy of sequential decision environments by domain and dynamics model. \dreamprice{} is the first learned world model for economic domains.}
\label{tab:gap-matrix}
\begin{tabular}{@{}lll@{}}
\toprule
& \textbf{Learned model} & \textbf{Hand-crafted model} \\
\midrule
\textbf{Physical domain} & DreamerV1--V3 \citep{hafner2020dreamerv1, hafner2021dreamerv2, hafner2023dreamerv3}, IRIS \citep{zhang2023iris}, TD-MPC2 \citep{hansen2024tdmpc2} & MuJoCo, Isaac Gym, CARLA \\
\textbf{Economic domain} & \textbf{\dreamprice{} (this work)} & ABIDES, DSGE models, AI Economist \citep{zheng2020aieconomist} \\
\bottomrule
\end{tabular}
\end{table}

Why do world models matter for pricing? Imagination-based planning enables counterfactual reasoning at scale. A retailer can ask: what if the price of SKU 3 were cut by 5\% next quarter? A learned world model rolls out latent trajectories under that hypothetical action, predicting demand responses, substitution effects, and downstream margin impact without executing the change in the real world. Such queries are infeasible with model-free methods, which require either costly online exploration or restrictive assumptions about the policy class. Offline learning from historical scanner data \citep{dominicks1999} further aligns with the constraints of retail: experimentation is expensive, and the Dominick's dataset alone spans 400 weeks across 93 stores and roughly 18,000 UPCs. The structured economic state space---prices, quantities, demographics, and derived features in 100--300 dimensions---is fundamentally more tractable than the 64$\times$64$\times$3 image observations that DreamerV3 was designed for \citep{hafner2023dreamerv3}. A world model trained on such tabular data can capture demand dynamics, promotional effects, and competitive structure without the burden of high-dimensional visual encoding.

Economic data, however, poses a challenge absent in physical domains: endogeneity. In physical simulators, the transition dynamics are exogenous to the agent; the laws of physics do not depend on the policy. In retail, prices and demand are simultaneously determined. Retailers set prices in response to expected demand, costs, and competitive pressure; demand, in turn, responds to those prices. Naive machine learning that regresses demand on price learns a confounded relationship, conflating the causal effect of price with selection and omitted variables \citep{hausman1978specification}. The resulting elasticity estimates are biased, and a world model built on such estimates would produce misleading counterfactuals. The econometric literature has long addressed this through instrumental variables \citep{hausman1978specification}, control functions \citep{berry1995automobile}, and more recently double machine learning \citep{chernozhukov2018dml, bach2024doubleml}. The Dominick's panel structure, with 93 stores and cross-store price variation, supports Hausman-style instruments: the average log price in other stores serves as an instrument for a given store's price, as it is relevant (prices are correlated across stores) but plausibly exogenous to store-specific demand shocks. \dreamprice{} integrates this identification strategy by pre-estimating causal price elasticities via DML-PLIV and freezing them into a constrained demand decoder. The decoder's residual MLP learns everything else---seasonality, demographics, promotion effects---while the price-demand relationship remains causally identified.

The contributions of this paper are as follows:

\begin{enumerate}[leftmargin=*]
\item \textbf{First learned world model for economic environments.} \dreamprice{} is the first system to learn a dynamics model for retail pricing from scanner data, enabling imagination-based planning and counterfactual demand estimation in a domain previously served only by hand-crafted simulators or model-free methods.

\item \textbf{Mamba-2 SSM backbone.} The recurrent backbone of standard RSSMs is replaced with Mamba-2 \citep{dao2024mamba2}, a selective state-space model that achieves $O(n)$ parallel training via structured state-space duality (SSD) and $O(1)$-per-step recurrent rollouts during imagination. This replaces the GRU-based recurrence of \dreamer{} \citep{hafner2023dreamerv3} with a more scalable alternative grounded in the S4 \citep{gu2022s4} and Mamba \citep{gu2024mamba} lineage.

\item \textbf{DRAMA-style decoupled posterior.} Following \citet{rajbhandari2024drama}, the posterior over the stochastic latent depends only on the current observation, not on the recurrent hidden state. This decoupling breaks the sequential dependency that would otherwise prevent Mamba-2's parallel SSD scan during training. During imagination, the model switches to recurrent \texttt{step()} mode for single-step rollouts.

\item \textbf{Causally-constrained demand decoder.} Per-category price elasticities are pre-estimated via DML-PLIV \citep{chernozhukov2018dml} and frozen into the \dreamprice{} decoder. The decoder enforces the causal price-demand relationship while an MLP residual learns seasonality, demographics, and promotion effects. This prevents the world model from learning confounded elasticities from the observational data.

\item \textbf{MOPO-style offline pessimism.} To guard against distributional shift inherent in offline reinforcement learning \citep{levine2020offline}, \dreamprice{} employs a 5-head reward ensemble and uses the lower confidence bound $r_{\text{pessimistic}} = r_{\text{mean}} - \lambda_{\text{LCB}} \cdot r_{\text{std}}$ for policy optimization \citep{yu2020mopo}. The actor is trained on this pessimistic reward signal while the critic and world model remain unchanged.
\end{enumerate}

The remainder of this paper is organized as follows. Section~\ref{sec:literature} reviews related work on world models, state-space models, causal demand estimation, and offline reinforcement learning. Section~\ref{sec:methodology} presents the \dreamprice{} architecture, training procedure, and causal identification strategy. Section~\ref{sec:results} reports experiments on the Dominick's canned soup category, including world model quality, elasticity recovery, and offline policy performance. Section~\ref{sec:conclusion} concludes and discusses limitations and future directions.
